\section{Conclusion and Future Work}

\subsection{Conclusion}
This project successfully transforms an RTOS starter code base into a complete
ESP32-S3 IoT system that combines sensing, local feedback, web interaction,
TinyML inference, and cloud telemetry. From the perspective of assignment
coverage, all six required tasks are addressed in a coherent way rather than as
independent features. Temperature is mapped to LED blink patterns, humidity is
mapped to NeoPixel color, the LCD provides multiple display states, the AP web
server supports both monitoring and control, the TinyML model runs on-device,
and CoreIOT publishing is handled correctly in STA mode.

From the perspective of system design, the most important achievement is the
task-based architecture. The application is structured around clear RTOS
responsibilities, explicit queue-based communication, and proper use of
semaphores and mutexes. This makes the system easier to understand, extend, and
test. The introduction of an application context and protected shared state also
improves code quality by reducing uncontrolled global-variable usage.

The project also provides meaningful engineering lessons. It shows how a single
embedded node can support multiple layers of functionality at once: physical
sensing, local actuation, browser-based configuration, machine-learning
inference, and cloud communication. It also highlights the importance of making
tradeoffs explicit. For example, the system favors batching for TinyML
stability, latest-value queues for simplicity, and AP/STA mode switching for
practical usability.

\subsection{What Was Learned}
Through the development of this system, the following technical lessons were
especially important:
\begin{itemize}
  \item how to decompose a real embedded application into RTOS tasks with
  different priorities and activation models;
  \item how to use queues, semaphores, and mutexes for their proper purposes
  rather than treating them as interchangeable mechanisms;
  \item how to design a user-facing AP dashboard that remains connected to live
  firmware state;
  \item how to integrate TinyML into a streaming sensor pipeline on a resource-
  constrained microcontroller;
  \item how cloud connectivity changes the state machine of an embedded node,
  especially when AP and STA modes must coexist.
\end{itemize}

\subsection{Future Work}
Although the implemented system already meets the project goals, there are
several clear paths for improvement:
\begin{itemize}
  \item improve long-term timing analysis with explicit task profiling and
  deadline measurements;
  \item integrate NTP synchronization directly into the STA connection path;
  \item evaluate the TinyML model with a larger on-device test set and report
  formal accuracy metrics;
  \item extend CoreIOT interaction with additional RPC-controlled features;
  \item store longer-term sensor history and expose it through the web
  interface;
  \item add recovery features such as offline telemetry buffering during
  temporary Wi-Fi loss.
\end{itemize}

Overall, the implemented system shows that an ESP32-S3 platform can host a
well-structured and feature-rich IoT application when its software is designed
around explicit RTOS principles rather than ad-hoc control logic.
